\documentclass{../note}

\begin{document}

\begin{zadanie}
Niech $f:\mathbb{R}\rightarrow\mathbb{R}$ będzie funkcją ciągłą. Czy funkcja może przyjmować każdą wartość rzeczywistą dokładnie 3 razy? A dokładnie 2 razy?
\end{zadanie}
\begin{rozwiazanie}
Niech
\[f(x) = \begin{cases}
    \lfloor\frac{x}{3}\rfloor + \{x\} \text{ gdy } \lfloor x\rfloor \equiv 0 \mod 3\\
    \lfloor\frac{x}{3}\rfloor + 1 + \{x\} \text{ gdy } \lfloor x\rfloor \equiv 1 \mod 3\\
    \lfloor\frac{x}{3}\rfloor + 2 - \{x\} \text{ gdy } \lfloor x\rfloor \equiv 1 \mod 3\\
\end{cases}\]
$f$ składa się z kawałków liniowych, więc żeby sprawdzić ciągłość, wystarczy sprawdzić ciągłość w punktach, w których zmienia się przypadek, czyli dla $x \in \Z$. Gdy $x \equiv 0 \mod 3$, to jest to pierwszy przypadek i wartość wynosi $\frac{x}{3}$, ale to też jest granica trzeciego przypadku:
\[\lim_{y\to x^-} f(y) = \lim_{y\to x^-} \lfloor\frac{y}{3}\rfloor + 2 - \{y\} = \lim_{y\to x^-} \frac{x}{3} - 1 + 2 - \{y\} = \frac{x}{3} - 1 + 2 - 1 = \frac{x}{3}\]
Nie będę rozpisywał wszystkich przypadków, ale to jest oczywiste gdy się spojrzy na wykres, bo to taki zygzak który idzie w górę. Zauważmy też, że każdy z przypadków przyjmuje każdą wartość dokładnie raz, w sensie że np. dla każdego $y$ istnieje dokładnie jedno $x$, takie że $\lfloor\frac{x}{3}\rfloor + \{x\} = y$, a mianowicie $x = 3\lfloor y \rfloor + \{y\}$ (bo ten pierwszy przypadek to jakby wziąć funkcję $y = x$, pociąć na kawałki i rozłożyć te kawałki w odstępach, żeby pokrywały dziedzinę tego przypadku i podobnie dla pozostałych dwóch przypadków). Tak więc każda liczba rzeczywista jest wartością $f$ dokładnie 3 razy.\\
Gdyby jakieś $f$ przyjmowało każdą wartość 2 razy, to istniałyby takie $a$ i $b$, że $f(a) = f(b) = 0$. Jeśli $f$ jest stałe na $[a, b]$, to $f$ przyjmuje wartość 0 nieskończenie wiele razy, co jest sprzecznością, więc $f$ musi się jakoś wychylać na tym przedziale, czyli musi być albo dodatnie supremum, albo ujemne infimum, to BSO niech to będzie że $s = \sup_{x\in[a, b]} f(x)$ i niech $c \in [a, b]$ będzie takie, że $f(c) = s$ (jest w skrypcie, że ono istnieje). $f$ nie przyjmuje wartości $s + 1$ na $[a, b]$, więc musi przyjmować tę wartość albo w $(-\infty, a)$, albo w $(b, \infty)$, to BSO niech to będzie, że $d \in (b, \infty)$ i $f(d) = s + 1$. Z tw. o wartościach pośrednich $f$ musi przyjąć wartość $\frac{s}{2}$ na $[b, d]$ (bo przyjmuje przedział wartości $[0, s + 1]$). Z tego samego twierdzenia $f$ musi też przyjmować wartość $\frac{s}{2}$ na $[a, c]$ i na $[c, a]$ (bo przyjmuje wszystkie wartości w $[0, s]$ na tych przedziałach). Zatem wartość $\frac{s}{2}$ jest przyjęta przez $f$ przynajmniej trzykrotnie, co jest sprzecznością.
\end{rozwiazanie}

\begin{zadanie}{3}
Niech $f:[0,1]\rightarrow\mathbb{R}$ będzie funkcją ciągłą spełniającą $f(0)=f(1).$ Wykaż, że istnieje $t\in(0,1)$, dla którego $f(t)=f(t^{2})$.
\end{zadanie}
\begin{rozwiazanie}
Jeśli $f$ jest stałe, to teza jest oczywista, a tak to musi mieć wartości większe od $f(0)$ albo mniejsze, to BSO niech będzie, że większe. Niech $s = \sup_{t\in(0, 1)} f(t)$ i niech $x\in(0, 1)$ jest takie, że $f(x) = s$ (jest w skrypcie, że ono istnieje). Niech $g(t) = f(t) - f(t^2)$. Zauważmy, że $t^2 \in (0,1)$, więc $g$ jest zdefiniowane na $(0, 1)$. Jest też ciągłe na tym przedziale, bo to różnica złożeń funkcji ciągłych. Skoro supremum $f$ jest w $x$, to $g(x) = f(x) - f(x^2) \geq 0$ i $g(\sqrt x) = f(\sqrt x) - f(x) \leq 0$, więc z tw. o wartościach pośrednich $g$ musi przyjąć wartość 0 na przedziale $[x, \sqrt x]$, czyli istnieje takie $t$, że $f(t) = f(t^2)$.
\end{rozwiazanie}

\end{document}